\documentclass[10pt,a4paper]{amsart}
\usepackage[margin=19mm]{geometry}
\usepackage{amsmath,amssymb,mathtools}
\usepackage[scr=rsfs,cal=euler]{mathalfa}
\usepackage{xfrac}
\usepackage[unicode,hidelinks,bookmarks=false]{hyperref}
\newcommand{\ie}{i.e.\ }
\newcommand{\cf}{cf.\ }
\newcommand{\resp}{respectively\ }
\newcommand{\Subsection}{\subsection}
\providecommand{\nobreakdash}{\nobreak\hbox{-}}
\setlength{\emergencystretch}{2em}
\multlinegap0pt
\usepackage{amssymb}
\usepackage{mathtools}
\usepackage[shortlabels]{enumitem}
\usepackage{xcolor}
\newtheorem{lemma}{Lemma}
\newcommand{\R}{\mathbb R}
\newcommand{\Z}{\mathbb Z}
\newcommand{\diam}{\operatorname{diam}}
\newcommand{\eps}{\varepsilon}
\title{The uniform Littlewood conjecture fails on a set of positive Hausdorff dimension}
\author{Nikita Shulga}
\date{Source: arXiv:2608.25059v1 (CC BY 4.0)}
\begin{document}
\maketitle

\noindent\emph{Source and licence.} The uniform Littlewood conjecture fails on a set of positive Hausdorff dimension. Version arXiv:2608.25059v1 (CC BY 4.0), distributed by arXiv under the Creative Commons Attribution 4.0 International licence, http://creativecommons.org/licenses/by/4.0/, as recorded in the arXiv licence metadata for this exact version and shown on the abstract page https://arxiv.org/abs/2608.25059v1. Full licence notice: \texttt{licenses/arxiv-2608.25059-CC-BY-4.0.txt}.\par\smallskip

\noindent\textit{Editorial scope.} Counted unit: the article's lemma lem:pair-rectangles (source0.tex lines 781--808). The article's complete original proof of it is delivered with it (source0.tex lines 810--1071, one \texttt{proof} environment of 262 lines). No mathematics is written, repaired or re-derived: the statement and the proof are the article's own lines, in the article's own notation, and no retained line is substituted or re-flowed.  Retained uncounted as the source-local context the proof needs: lines 449--456; lines 474--483; lines 556--564.  Nothing was omitted from any retained span, so the package carries no omission record.  No retained line contains a citation, so the wrapper carries no bibliography and no bibliography entry.  No retained line contains a figure environment, an image-inclusion command or drawing code, so no image is delivered and the package declares no asset dependency.  Editorial formatting only, in addition: an AMS \texttt{amsart} preamble that restates the article's own declarations and macros used by the retained lines, namely the environments \texttt{lemma} on separate counters, together with the standard packages those lines need.  The retained environments print in this document as Lemma 1 in order of appearance, whereas the article prints the same items with the numbering of its own full document.  The complete native source of the article as distributed in the arXiv e-print for this exact version is bundled unchanged as \texttt{sources/208541/source0.tex}.
\begin{lemma}\label{lem:local-rational-upper}
Let $A\ge2$, let $\mathbf{u}$ be an admissible word, and let $U\subset\R$ be a bounded
interval.  Uniformly in $\mathbf{u}$, $U$, and $T\ge2$,
\begin{equation}\label{eq:hensley-local-upper}
 \sum_{T\le q<6T}|M_A(q;F_A(\mathbf{u})\cap U)|
 \ll_A (1+|U|T^2)^{\delta_A}.
\end{equation}
\end{lemma}

\begin{lemma}\label{lem:fixed-prefix-count}
Let $A\ge2$ and let $\mathbf{u}$ be an admissible word.  There is
$T_0=T_0(\mathbf{u})$ such that, for every $T\ge T_0$,
\begin{equation}\label{eq:fixed-prefix-count}
 \sum_{T\le q<6T}|M_A(q;F_A(\mathbf{u}))|
 \asymp_A q(\mathbf{u})^{-2\delta_A}T^{2\delta_A}
 \asymp_A (\diam F_A(\mathbf{u}))^{\delta_A}T^{2\delta_A}.
\end{equation}
The implicit constants are independent of $\mathbf{u}$.
\end{lemma}

\begin{lemma}\label{lem:bilinear-count}
Let $q\ge2$, let $P,M\subseteq(\Z/q\Z)^\times$, and let
$H\subset\Z/q\Z$ be an interval of length $h$.  For every $\eps>0$,
\begin{equation}\label{eq:bilinear-count}
 \#\{(p,m)\in P\times M:pm\in H\}
 =\frac{h}{q}|P||M|
 +O_\eps\left(q^{1/2+\eps}(|P||M|)^{1/2}\right).
\end{equation}
\end{lemma}

\begin{lemma}\label{lem:pair-rectangles}
Assume
\begin{equation}\label{eq:beta-range}
 \max\left\{2-2\delta,\,-1+\sqrt{5-2\delta}\right\}<\beta<2\delta-1,
\end{equation}
and put
\begin{equation}\label{eq:s-a-def}
 s=\delta+\frac\beta2,\qquad \alpha=\delta+1.
\end{equation}
There is $\chi>0$ such that, for every admissible word $\mathbf{u}$ and every
nondegenerate interval $J\subset(0,1)$, there is $T_0=T_0(\mathbf{u},J)$ for
which every $T\ge T_0$ admits a finite family $\mathcal R_T(\mathbf{u},J)$ of
admissible $T$-rectangles satisfying
\begin{equation}\label{eq:pair-rectangle-lower}
 \#\mathcal R_T(\mathbf{u},J)
 \gg (\diam F_A(\mathbf{u}))^\delta |J|T^{2s}.
\end{equation}
The rectangles are pairwise separated by $\gg T^{-2}$.  Moreover, every ball
$B(z,\rho)\subset\R^2$ satisfies
\begin{align}
 \#\{Q\in\mathcal R_T(\mathbf{u},J):Q\cap B(z,\rho)\ne\varnothing\}
 &\ll (1+\rho T^2)^s &&(\rho\ge T^{-2}),
 \label{eq:pair-local-s}\\
 \#\{Q\in\mathcal R_T(\mathbf{u},J):Q\cap B(z,\rho)\ne\varnothing\}
 &\ll \rho^\alpha T^{2s} &&(\rho\ge T^{-\chi}).
 \label{eq:pair-local-a}
\end{align}
\end{lemma}

\begin{proof}
Choose $\eps>0$ so small that
\begin{equation}\label{eq:eps-choice-pair}
 s>1+\eps,
 \qquad
 \left(1-\frac\beta2\right)
 \frac{2+\beta/2+\eps}{1+\delta/2}<1.
\end{equation}
The first inequality is possible because $s>1$, which is equivalent to
$\beta>2-2\delta$. The second inequality is needed for some local estimate in Step~3 below. 

Fix
\begin{equation}\label{eq:chi-choice}
 0<\chi<\min\left\{1,
 \frac{s-1-\eps}{1+\delta/2}\right\}.
\end{equation}

\medskip
\noindent\emph{Step 1: construction and total count.}
Call $q\in[T,6T)$ \emph{rich} if $|M_A(q)|\ge T^\beta$.  For each
rich $q$, choose a subset $M_q'\subset M_A(q)$ with
$|M_q'|\asymp T^\beta$. For instance, for large $T$ we may require
$T^\beta/2\le |M_q'|\le T^\beta$.  Consider the total prefix count
$$
 S(T;\mathbf{u})=\sum_{T\le q<6T}|M_A(q;F_A(\mathbf{u}))|,
$$
which counts pairs $(q,p)$ with $T\le q<6T$ and $p\in M_A(q;F_A(\mathbf{u}))$.
If $q$ is not rich, then
$|M_A(q;F_A(\mathbf{u}))|\le |M_A(q)|<T^\beta$.  Since there are $O(T)$ possible
moduli $q\in[T,6T)$, the total contribution of the non-rich moduli to
$S(T;\mathbf{u})$ is at most
$$
 \sum_{\substack{T\le q<6T\\ q\ \mathrm{not\ rich}}}
 |M_A(q;F_A(\mathbf{u}))|\ll T^{1+\beta}.
$$
The inequality $1+\beta<2\delta$ and Lemma~\ref{lem:fixed-prefix-count} show that
$S(T;\mathbf{u})\gg (\diam F_A(\mathbf{u}))^\delta T^{2\delta}$.  Hence the rich moduli alone
already contribute
\begin{equation}\label{eq:high-prefix-mass}
 \sum_{\substack{T\le q<6T\\q\ \mathrm{rich}}}
 |M_A(q;F_A(\mathbf{u}))|
 \gg (\diam F_A(\mathbf{u}))^\delta T^{2\delta}.
\end{equation}

Choose an interval $J'\Subset J$ with $|J'|\ge |J|/3$.  For sufficiently
large $T$, every interval of radius $K_Aq^{-2}$ centered at a point
$r/q\in J'$ is contained in $J$.  Let
$$
 H_q=\{r\bmod q:1\le r<q,\ r/q\in J'\}.
$$
Then $|H_q|\gg |J|q$.  For a rich modulus $q$, let
$$
 N_q:=\#\{(p,b)\in M_A(q;F_A(\mathbf{u}))\times M_q': pb\bmod q\in H_q\}.
$$
For every pair counted by $N_q$, let $r\in\{1,\ldots,q-1\}$ be the least
positive residue satisfying $r\equiv pb\pmod q$.  Since $r/q\in J'$, the corresponding set
$$
 Q(q,p,b):=F_A(\mathbf{w}(p/q))\times
 \left[\frac rq-\frac{K_A}{q^2},\frac rq+\frac{K_A}{q^2}\right]
$$
is an admissible rectangle contained in $F_A(\mathbf{u})\times J$.  Let
$\widetilde{\mathcal R}_T(\mathbf{u},J)$ be the family of all rectangles $Q(q,p,b)$
obtained in this way.  Applying Lemma~\ref{lem:bilinear-count} with
$$
 P=M_A(q;F_A(\mathbf{u})),\qquad M=M_q',\qquad H=H_q
$$
therefore gives
$$
 N_q=\frac{|H_q|}{q}\,|M_A(q;F_A(\mathbf{u}))|\,|M_q'|
 +O_\eps\!\left(q^{1/2+\eps}|M_A(q;F_A(\mathbf{u}))|^{1/2}T^{\beta/2}\right).
$$
Since $|H_q|\gg |J|q$ and $|M_q'|\asymp T^\beta$, the main term for a
fixed rich modulus is
$$
 \gg |J|T^\beta|M_A(q;F_A(\mathbf{u}))|.
$$
Summing this lower bound over all rich moduli and using
\eqref{eq:high-prefix-mass}, we obtain a total main-term contribution
$$
 \gg |J|T^\beta\sum_{\substack{T\le q<6T\\ q\ \mathrm{rich}}}|M_A(q;F_A(\mathbf{u}))|
 \gg (\diam F_A(\mathbf{u}))^\delta |J|T^{2\delta+\beta}
 =(\diam F_A(\mathbf{u}))^\delta |J|T^{2s}.
$$
The sum of the error terms is, by Cauchy--Schwarz and the upper bound in
Lemma~\ref{lem:fixed-prefix-count},
\begin{align*}
 T^{1/2+\beta/2+\eps}
 \sum_{T\le q<6T}|M_A(q;F_A(\mathbf{u}))|^{1/2}
 &\ll T^{1+\beta/2+\eps}
 \left(\sum_{T\le q<6T}|M_A(q;F_A(\mathbf{u}))|\right)^{1/2}\\
 &\ll T^{\delta+1+\beta/2+\eps}.
\end{align*}
Since $s>1+\eps$, the error is $o(T^{2s})$.  For each fixed
$(\mathbf{u},J)$, the factor $(\diam F_A(\mathbf{u}))^\delta|J|$ in the
main term is positive, so $T_0(\mathbf{u},J)$ may be enlarged until the error
is at most half of that main term.  Thus the lower-bound constant below is
independent of $\mathbf{u}$ and $J$, although the threshold $T_0(\mathbf{u},J)$
may depend on them:
$$
 \sum_{\substack{T\le q<6T\\q\ \mathrm{rich}}}N_q
 \gg (\diam F_A(\mathbf{u}))^\delta |J|T^{2s}.
$$

The map from counted triples $(q,p,b)$ to the associated points $(p/q,r/q)$ is injective.
Suppose first that $p/q\ne p'/q'$.  Both fractions are reduced and have
denominators in $[T,6T)$, so
$$
 \left|\frac pq-\frac{p'}{q'}\right|\ge\frac1{qq'}\ge\frac1{36T^2}.
$$
If $p/q=p'/q'$, reducedness gives $p=p'$ and $q=q'$.  Two distinct counted
triples must then have $b\ne b'$, and multiplication by $p$ is invertible modulo
$q$, so the corresponding residues $r\equiv pb\pmod q$ and
$r'\equiv pb'\pmod q$ are distinct.  Their second coordinates differ by
at least $q^{-1}$.  Thus the map $(q,p,b)\mapsto(p/q,r/q)$ is injective.  Moreover, distinct
canonical fractions have distinct canonical words and hence distinct
first-coordinate cylinders.  Therefore distinct counted triples yield distinct
rectangles, and their associated points are $\gg T^{-2}$-separated.
In particular,
$$
 \#\widetilde{\mathcal R}_T(\mathbf{u},J)
 =\sum_{\substack{T\le q<6T\\q\ \mathrm{rich}}}N_q
 \gg (\diam F_A(\mathbf{u}))^\delta |J|T^{2s}.
$$
We also have
\begin{equation}\label{eq:pair-total-upper}
 \#\widetilde{\mathcal R}_T(\mathbf{u},J)
 \le T^\beta\sum_{T\le q<6T}|M_A(q;F_A(\mathbf{u}))|
 \ll T^{2s}.
\end{equation}

For $Q=Q(q,p,b)\in\widetilde{\mathcal R}_T(\mathbf{u},J)$, write
$c(Q)=(p/q,r/q)$.  The preceding argument gives
$|c(Q)-c(Q')|\ge (36T^2)^{-1}$ for distinct associated points, while every
rectangle has diameter $O_A(T^{-2})$.  Choose a maximal subfamily whose
associated points are separated by $C_AT^{-2}$, with $C_A$ larger than twice
the rectangle-diameter constant.  The $(36T^2)^{-1}$ separation shows that
each selected point excludes only $O_A(1)$ members, so the subfamily has
cardinality comparable to $\widetilde{\mathcal R}_T(\mathbf{u},J)$ and its
rectangles are mutually $\gg_A T^{-2}$-separated.  Thus
\eqref{eq:pair-rectangle-lower} is preserved.

\medskip
\noindent\emph{Step 2: local counting estimates.}
It is enough to estimate the larger family
$\widetilde{\mathcal R}_T(\mathbf{u},J)$, since $\mathcal R_T(\mathbf{u},J)$ is a subfamily.
Fix a ball
$B(z,\rho)$ with $\rho\ge T^{-2}$, let $N(\rho)$ be the number of rectangles
meeting it, and put
$$
 X=1+\rho T^2.
$$
If a rectangle meets the ball, its associated point $(p/q,r/q)$ lies in a product $I_1\times I_2$
of intervals of length $O_A(\rho)$.  Define
$$
 P_q=\{p\in M_A(q;F_A(\mathbf{u})):p/q\in I_1\}.
$$
Lemma~\ref{lem:local-rational-upper} gives
\begin{equation}\label{eq:local-p-mass}
 \sum_{T\le q<6T}|P_q|\ll X^\delta.
\end{equation}

For fixed $q$ and $p\in P_q$, the map $b\mapsto pb\pmod q$ is injective, and
$I_2$ contains $O(1+\rho q)$ residue points of the form $r/q$.  Hence
\begin{equation}\label{eq:pair-crude-local}
 N(\rho)\ll (1+\rho T)X^\delta.
\end{equation}

For a second estimate, let
$H_{q,2}=\{r\bmod q:0\le r<q,\ r/q\in I_2\}$.  Its length is $O(1+\rho q)$.
For rich $q$, Lemma~\ref{lem:bilinear-count}, applied with $P=P_q$,
$M=M_q'$, and $H=H_{q,2}$, gives after summing over the rich moduli
\begin{align}
 N(\rho)
 &\ll (\rho+T^{-1})T^\beta\sum_{\substack{T\le q<6T\\q\ \mathrm{rich}}}|P_q|
 +T^{1/2+\beta/2+\eps}\sum_{\substack{T\le q<6T\\q\ \mathrm{rich}}}|P_q|^{1/2}\notag\\
 &\ll (\rho+T^{-1})T^\beta X^\delta
 +T^{1+\beta/2+\eps}X^{\delta/2}.
 \label{eq:pair-bilinear-local}
\end{align}
In the last line we used \eqref{eq:local-p-mass} and Cauchy--Schwarz.

\medskip
\noindent\emph{Step 3: proof of \eqref{eq:pair-local-s}.}
If $X\le2T$, then $1+\rho T\ll1$, and \eqref{eq:pair-crude-local} gives
$N(\rho)\ll X^\delta\le X^s$.  If $X\ge T^2$, \eqref{eq:pair-total-upper} gives
$N(\rho)\ll T^{2s}\le X^s$.

It remains to consider $2T<X<T^2$.  In this range
$\rho+T^{-1}\ll X/T^2$, so the main term in
\eqref{eq:pair-bilinear-local} is at most
$$
 T^{\beta-2}X^{\delta+1}
 =X^s\left(\frac{X}{T^2}\right)^{1-\beta/2}
 \le X^s.
$$
For the other contribution, use the better of the crude bound and the error
term in \eqref{eq:pair-bilinear-local}.  After division by $X^s$, these two
bounds become
$$
 f_1(X)=T^{-1}X^{1-\beta/2},
 \qquad
 f_2(X)=T^{1+\beta/2+\eps}X^{-(\delta+\beta)/2}.
$$
The function $f_1$ increases and $f_2$ decreases.  At the left endpoint,
$$
 f_1(2T)\ll T^{-\beta/2}=o(1),
 \qquad
 f_2(2T)\asymp T^{1+\eps-\delta/2}\to\infty,
$$
whereas at the right endpoint,
$$
 f_1(T^2)=T^{1-\beta}\to\infty,
 \qquad
 f_2(T^2)=T^{1+\eps-s}=o(1).
$$
Here $0<\beta<1$ and $s>1+\eps$.  Thus the graphs cross exactly once in the
intermediate range.  At the crossing,
$$
 X=T^{(2+\beta/2+\eps)/(1+\delta/2)},
$$
and their common value is
$$
 T^{-1+(1-\beta/2)(2+\beta/2+\eps)/(1+\delta/2)}.
$$
Thus, in order that the better of the two bounds remain uniformly bounded
(and in fact be $o(1)$ at the crossing), we need
$$
 \left(1-\frac\beta2\right)
 \frac{2+\beta/2+\eps}{1+\delta/2}<1,
$$
which is the second condition in \eqref{eq:eps-choice-pair}.  Setting
$\eps=0$ and expanding gives
$$
 \beta^2+2\beta+2\delta-4>0.
$$
Since $\beta>0$, this is equivalent to
$$
 \beta>-1+\sqrt{5-2\delta}.
$$
This is the origin of the extra lower bound on $\beta$ in \eqref{eq:beta-range}.
Since $\min\{f_1,f_2\}$ increases up to the crossing and decreases afterwards,
it is bounded throughout the intermediate range.  This proves
\eqref{eq:pair-local-s}.

\medskip
\noindent\emph{Step 4: proof of \eqref{eq:pair-local-a}.}
For $\rho\ge1$ the claim follows from
the total bound \eqref{eq:pair-total-upper}, since $\rho^\alpha\ge1$.
Thus assume $T^{-\chi}\le\rho\le1$.  Since $\chi<1$, we have
$\rho\gg T^{-1}$ and $X\asymp\rho T^2$.  The main term in
\eqref{eq:pair-bilinear-local} is then
$$
 \ll \rho T^\beta(\rho T^2)^\delta
 =\rho^\alpha T^{2s}.
$$
The ratio of the error term to $\rho^\alpha T^{2s}$ is
$$
 T^{1-s+\eps}\rho^{-1-\delta/2}
 \le T^{1-s+\eps+\chi(1+\delta/2)}\le1
$$
by \eqref{eq:chi-choice}.  This completes the proof.
\end{proof}

\end{document}
